%============================================================
% CH2 · P1 · PROBLEM 1 — ROUND TRIP (CSES 1669)
%============================================================
\subsection{Problem 2.1.1: Round Trip}
\problemlabel{2.1.1}{ssec:ch2-round-trip}

\problemcard{2.1.1}{Round Trip}{CSES 1669}{https://cses.fi/problemset/task/1669}{\exeasy\ (CSES Graph Algorithms; the classic first cycle problem)}{Undirected cycle detection: DFS, skip the parent, a visited neighbour closes a cycle; print it from the parent pointers}

\subsubsection*{1. Problem statement}

Byteland has $n$ cities and $m$ two-way roads. Design a \textbf{round trip} that starts in a city, visits two or more other cities and returns to the start; all intermediate cities must be distinct. Print the number $k$ of cities on the route (the start counted twice) and the route itself, or \texttt{IMPOSSIBLE}. Any valid route is accepted.

\textbf{Constraints:} $1 \le n \le 10^5$, $1 \le m \le 2 \cdot 10^5$; every road joins two different cities, and there is at most one road between two cities (a \emph{simple} graph).

\begin{table}[H]
\centering\small
\begin{tabular}{@{}lll@{}}
\toprule
\textbf{Input} & \textbf{Official output} & \textbf{Our output} \\
\midrule
\texttt{5 6 / 1 3 / 1 2 / 5 3 / 1 5 / 2 4 / 4 5} & \texttt{4 / 3 5 1 3} & \texttt{4 / 1 5 3 1} \\
\bottomrule
\end{tabular}
\end{table}

\begin{diagrambox}
\begin{tikzpicture}[x=1.15cm, y=1cm]
    \begin{scope}
        \node[gv] (a1) at (1.6,1.3) {1}; \node[gv] (a3) at (0,0.3) {3}; \node[gv] (a2) at (3.2,0.3) {2};
        \node[gv] (a5) at (1.6,-0.7) {5}; \node[gv] (a4) at (3.2,-1.2) {4};
        \foreach \a/\b in {1/3, 1/2, 3/5, 1/5, 2/4, 4/5} \draw[ge] (a\a) -- (a\b);
        \node[font=\scriptsize\bfseries] at (1.6,-1.9) {the sample: 5 cities, 6 roads};
    \end{scope}
    \begin{scope}[xshift=5.4cm]
        \node[gvsel] (b1) at (1.6,1.3) {1}; \node[gvsel] (b3) at (0,0.3) {3}; \node[gv] (b2) at (3.2,0.3) {2};
        \node[gvsel] (b5) at (1.6,-0.7) {5}; \node[gv] (b4) at (3.2,-1.2) {4};
        \foreach \a/\b in {1/2, 2/4, 4/5} \draw[ge] (b\a) -- (b\b);
        \foreach \a/\b in {1/3, 3/5, 1/5} \draw[gpath] (b\a) -- (b\b);
        \node[font=\scriptsize\bfseries] at (1.6,-1.9) {our answer: $1 \to 5 \to 3 \to 1$};
    \end{scope}
    \begin{scope}[xshift=10.8cm]
        \node[gvsel] (c1) at (1.6,1.3) {1}; \node[gv] (c3) at (0,0.3) {3}; \node[gvsel] (c2) at (3.2,0.3) {2};
        \node[gvsel] (c5) at (1.6,-0.7) {5}; \node[gvsel] (c4) at (3.2,-1.2) {4};
        \foreach \a/\b in {1/3, 3/5} \draw[ge] (c\a) -- (c\b);
        \foreach \a/\b in {1/2, 2/4, 4/5, 1/5} \draw[gpath] (c\a) -- (c\b);
        \node[font=\scriptsize\bfseries] at (1.6,-1.9) {also valid: $1, 2, 4, 5, 1$};
    \end{scope}
\end{tikzpicture}
\figcap{The sample has three cycles ($1,3,5$; $1,2,4,5$; $1,2,4,5,3$); the judge accepts any of them. The official answer $3,5,1,3$ is the first triangle read from 3.}
\end{diagrambox}

\subsubsection*{2. The graph model}

\begin{modelbox}[Graph Model --- Round Trip]
\renewcommand{\arraystretch}{1.35}
\begin{tabularx}{\linewidth}{@{}>{\bfseries}l X@{}}
Vertices & the $n$ cities \\
Edges & undirected, one per road; simple graph (no loops, no parallel edges); possibly disconnected \\
Weights & none \\
Question & find any cycle; a cycle in a simple graph automatically has $\ge 3$ vertices, which is exactly ``two or more other cities'' \\
Size & $|V| \le 10^5$, $|E| \le 2\cdot10^5$: we need \VE; a DFS path can be $10^5$ deep \\
\end{tabularx}
\end{modelbox}

\subsubsection*{3. Key observations}

\begin{itemize}[leftmargin=*]
    \item \textbf{``Round trip with distinct intermediate cities'' $=$ ``cycle''.} The route repeats only its start, so it is a cycle of the graph; conversely every cycle is a valid route.
    \item \textbf{Length $\ge 3$ is free.} A cycle of length 1 needs a loop, one of length 2 needs two parallel roads; the input has neither. So the only thing to avoid is the fake ``cycle'' $p, v, p$ made by walking a tree edge back --- skip the parent.
    \item \textbf{Existence is easy, printing needs the DFS.} A forest with $c$ components has $n - c$ edges, so a cycle exists iff $m > n - c$. To print one we need the moment a non-tree edge is found.
    \item \textbf{Depth.} A path $1 - 2 - \dots - 10^5$ plus the road $10^5 - 1$ is a valid input whose DFS goes $10^5$ deep: write the DFS iteratively.
\end{itemize}

\subsubsection*{4. Approaches}

\begin{bruteforcebox}
\textbf{Idea: an edge lies on a cycle iff its ends stay connected without it.} Try every road $(a,b)$: run a BFS from $a$ that is not allowed to use this road. If it reaches $b$, the BFS path $a \leadsto b$ plus the road is a cycle.

\begin{lstlisting}[style=cpp]
#include <bits/stdc++.h>
using namespace std;

// Brute force: an edge (a, b) lies on a cycle iff a and b stay connected without it.
// Try every edge; BFS from a to b avoiding that edge; the BFS path + the edge is the cycle.
int main() {
    int n, m;
    cin >> n >> m;
    vector<pair<int, int>> e(m);
    vector<vector<pair<int, int>>> adj(n + 1);          // (neighbour, edge id)
    for (int i = 0; i < m; i++) {
        cin >> e[i].first >> e[i].second;
        adj[e[i].first].push_back({e[i].second, i});
        adj[e[i].second].push_back({e[i].first, i});
    }
    for (int skip = 0; skip < m; skip++) {
        int a = e[skip].first, b = e[skip].second;
        vector<int> par(n + 1, -1);
        queue<int> q;
        q.push(a); par[a] = a;
        while (!q.empty()) {
            int u = q.front(); q.pop();
            for (auto& pr : adj[u])
                if (pr.second != skip && par[pr.first] == -1) { par[pr.first] = u; q.push(pr.first); }
        }
        if (par[b] == -1) continue;                     // edge skip is on no cycle
        vector<int> cyc;
        for (int x = b; x != a; x = par[x]) cyc.push_back(x);
        cyc.push_back(a);
        cyc.push_back(b);                               // close the loop with the skipped edge
        cout << cyc.size() << "\n";
        for (int x : cyc) cout << x << " ";
        cout << "\n";
        return 0;
    }
    cout << "IMPOSSIBLE\n";
}
\end{lstlisting}

On the sample the first road $(1,3)$ already works: the BFS from 1 avoiding it reaches 3 via 5, and the program prints \texttt{4 / 3 5 1 3} --- by chance exactly the official answer.

\begin{complexitybox}
\textbf{Time:} up to $m$ searches of \bigO{n + m} each, \bigO{m(n+m)}. A bad input: a tree on $10^5 - 3$ cities listed first, then a separate triangle --- every tree road costs a BFS over the whole tree before the triangle is reached, about $10^{10}$ steps. \textbf{Space:} \bigO{n+m}. \textbf{Verdict:} \tle{}, but a perfect reference for testing.
\end{complexitybox}
\end{bruteforcebox}

\begin{algorithmbox}[Approach 2: one DFS, the first non-tree edge closes the cycle]
\textbf{Idea.} Run DFS from every unvisited city. Sitting at $u$, skip the parent; any other neighbour $v$ that is already visited is an \textbf{ancestor} of $u$ (undirected DFS has only tree and back edges, Lesson~\ref{ssec:ch2-dfs-tree}), so $v, u, \texttt{par}[u], \dots, v$ is a cycle. Iterative, with \texttt{it[u]} as in Lesson~\ref{ssec:ch2-call-stack}.

\begin{lstlisting}[style=cpp]
#include <bits/stdc++.h>
using namespace std;

// Round Trip (CSES 1669): find a cycle of length >= 3 in a simple undirected graph.
// Iterative DFS; skipping the parent vertex is safe because the graph is simple.
int main() {
    ios::sync_with_stdio(false);
    cin.tie(nullptr);
    int n, m;
    cin >> n >> m;
    vector<vector<int>> adj(n + 1);
    for (int i = 0; i < m; i++) {
        int a, b;
        cin >> a >> b;
        adj[a].push_back(b);
        adj[b].push_back(a);
    }
    vector<int> par(n + 1, 0);
    vector<char> seen(n + 1, 0);
    vector<size_t> it(n + 1, 0);
    for (int s = 1; s <= n; s++) {
        if (seen[s]) continue;
        seen[s] = 1;
        vector<int> st(1, s);
        while (!st.empty()) {
            int u = st.back();
            if (it[u] == adj[u].size()) { st.pop_back(); continue; }   // u is finished
            int v = adj[u][it[u]++];
            if (v == par[u]) continue;                 // the tree edge we came in by
            if (!seen[v]) {                            // tree edge u - v
                seen[v] = 1; par[v] = u;
                st.push_back(v);
            } else {                                   // back edge u - v: v is an ancestor of u
                vector<int> cyc;
                cyc.push_back(v);
                for (int x = u; x != v; x = par[x]) cyc.push_back(x);
                cyc.push_back(v);                      // v, u, par[u], ..., v
                cout << cyc.size() << "\n";
                for (int x : cyc) cout << x << " ";
                cout << "\n";
                return 0;
            }
        }
    }
    cout << "IMPOSSIBLE\n";
}
\end{lstlisting}

\begin{complexitybox}
\textbf{Time:} \VE{} --- every city is pushed once, every road examined twice; the cycle is printed in \bigO{n}. \textbf{Space:} \bigO{n+m}, no recursion. \textbf{Verdict:} \ac{}; on a path of $10^5$ cities closed by one road it answers in well under a second (the printed cycle has all $10^5$ cities).
\end{complexitybox}
\end{algorithmbox}

\begin{invariantbox}[Why the visited neighbour is an ancestor]
\textbf{Invariant.} The stack holds the current path $x_0, \dots, x_k = u$ with $\texttt{par}[x_{i+1}] = x_i$; every finished vertex has had all its roads examined.

Suppose $u$ examines a road to a visited $v \ne \texttt{par}[u]$. If $v$ were finished and a descendant of $u$, then $v$ had examined the same road earlier, found $u$ visited and not its parent, and the program would already have stopped. A finished non-descendant is impossible in an undirected DFS: it would have entered $u$ itself while scanning that road. So $v$ is on the stack, an ancestor $x_j$ with $j \le k - 2$ ($j = k-1$ would be the parent, and there is only one road between two cities). The walk $v = x_j, x_k, x_{k-1}, \dots, x_{j+1}, x_j$ is a cycle of $k - j + 1 \ge 3$ cities.

If the search ends without stopping, every road was a tree edge, the graph is a forest, and \texttt{IMPOSSIBLE} is correct. \hfill$\blacksquare$
\end{invariantbox}

\begin{alternativebox}[Approach 3: DSU --- the first road that joins two connected cities]
Add the roads one by one to a disjoint-set union (Chapter~6). A road whose ends are already in the same set closes a cycle; the cycle is that road plus the unique path between its ends in the forest of the roads accepted so far, found by one BFS. Union by size keeps the sets shallow --- with the naive ``\texttt{dsu[find(a)] = find(b)}'' and a recursive \texttt{find}, the path input above builds a chain of $10^5$ parents, and our first version of this program crashed with a stack overflow inside \texttt{find}.

\begin{lstlisting}[style=cpp]
#include <bits/stdc++.h>
using namespace std;

// Alternative: add the roads one by one to a DSU. The first road (a, b) whose ends are already
// connected closes a cycle: the forest path from a to b plus the road itself.
vector<int> dsu, sz;
int findSet(int x) {                              // iterative, with path halving
    while (dsu[x] != x) x = dsu[x] = dsu[dsu[x]];
    return x;
}

int main() {
    ios::sync_with_stdio(false);
    cin.tie(nullptr);
    int n, m;
    cin >> n >> m;
    dsu.resize(n + 1); sz.assign(n + 1, 1);
    iota(dsu.begin(), dsu.end(), 0);
    vector<vector<int>> forest(n + 1);
    for (int i = 0; i < m; i++) {
        int a, b;
        cin >> a >> b;
        int ra = findSet(a), rb = findSet(b);
        if (ra != rb) {                                // a new forest edge: union by size
            if (sz[ra] < sz[rb]) swap(ra, rb);
            dsu[rb] = ra; sz[ra] += sz[rb];
            forest[a].push_back(b);
            forest[b].push_back(a);
            continue;
        }
        // a and b already connected: BFS in the forest for the unique path a ... b
        vector<int> par(n + 1, 0);
        par[a] = a;
        queue<int> q;
        q.push(a);
        while (!q.empty()) {
            int u = q.front(); q.pop();
            for (int v : forest[u])
                if (!par[v]) { par[v] = u; q.push(v); }
        }
        vector<int> cyc;
        for (int x = b; x != a; x = par[x]) cyc.push_back(x);
        cyc.push_back(a);
        cyc.push_back(b);
        cout << cyc.size() << "\n";
        for (int x : cyc) cout << x << " ";
        cout << "\n";
        return 0;
    }
    cout << "IMPOSSIBLE\n";
}
\end{lstlisting}
On the sample the roads $1{-}3$, $1{-}2$, $5{-}3$ are accepted, and $1{-}5$ is the first road inside one set: output \texttt{4 / 5 3 1 5}.
\begin{complexitybox}
\textbf{Time:} \bigO{m\,\alpha(n)} for the DSU plus one \bigO{n} BFS. \textbf{Space:} \bigO{n}. \textbf{Verdict:} \ac{}. The DSU version finds the cycle closed by the \emph{earliest} possible road in input order --- useful when the problem asks for that (``the first road that creates a loop'').
\end{complexitybox}
\end{alternativebox}

\begin{tricksbox}[Approach 4: Implementation Refinements]
\begin{itemize}[leftmargin=*]
    \item \textbf{Recursive version.} Ten lines shorter, but it recurses $10^5$ levels deep on the path input; on a 2\,MB stack (Lesson~\ref{ssec:ch2-call-stack}) that crashes. Unless you know the judge's stack is large, prefer the iterative form for $n \ge 5\cdot10^4$.
    \item \textbf{Parent vertex vs.\ parent edge.} Skipping \texttt{v == par[u]} is correct only because the graph is simple. For multigraphs store edge ids (template 1 of the pattern).
    \item \textbf{Fast output.} The cycle can have $10^5$ cities; print with one \texttt{"\textbackslash n"} at the end, never \texttt{endl} after every number.
    \item \textbf{Root's parent.} \texttt{par[s] = 0} works because cities are numbered from 1, so 0 is never a neighbour.
\end{itemize}
\end{tricksbox}

\subsubsection*{5. Trace}

\begin{tracebox}[{Sample: adjacency lists 1: 3 2 5,\ 2: 1 4,\ 3: 1 5,\ 4: 2 5,\ 5: 3 1 4}]
\begin{center}
\begin{tikzpicture}[x=0.85cm, y=0.85cm]
    \tikzset{gv/.append style={minimum size=6.5mm, font=\scriptsize\bfseries}}
    \foreach \k/\t/\sa/\sb/\sc/\sd/\se in {
        0/{1. enter 1; 1--3 is new}/gvcur/gv/gv/gv/gv,
        1/{2. at 3: skip 1; 3--5 is new}/gvseen/gv/gvcur/gv/gv,
        2/{3. at 5: skip 3; 5--1 is visited!}/gvsel/gv/gvsel/gv/gvcur}{
        \begin{scope}[xshift=\k*5.4cm]
            \node[\sa] (t1) at (1.6,1.3) {1}; \node[\sc] (t3) at (0,0.3) {3}; \node[\sb] (t2) at (3.2,0.3) {2};
            \node[\se] (t5) at (1.6,-0.7) {5}; \node[\sd] (t4) at (3.2,-1.2) {4};
            \foreach \a/\b in {1/2, 2/4, 4/5} \draw[ge] (t\a) -- (t\b);
            \ifnum\k=0 \draw[ge] (t1) -- (t3); \draw[ge] (t3) -- (t5); \draw[ge] (t1) -- (t5); \fi
            \ifnum\k=1 \draw[gtree] (t1) -- (t3); \draw[ge] (t3) -- (t5); \draw[ge] (t1) -- (t5); \fi
            \ifnum\k=2 \draw[gtree] (t1) -- (t3); \draw[gtree] (t3) -- (t5); \draw[gback, -] (t5) -- (t1); \fi
            \node[font=\scriptsize\bfseries] at (1.6,2.05) {\t};
        \end{scope}
    }
\end{tikzpicture}
\end{center}
\begin{center}\small
\begin{tabular}{@{}clll@{}}
\toprule
\textbf{step} & \textbf{at $u$} & \textbf{examine} & \textbf{action / stack after} \\
\midrule
1 & 1 & $1{-}3$: 3 unvisited & tree edge, \texttt{par[3] = 1}; stack \texttt{1 3} \\
2 & 3 & $3{-}1$: parent & skip \\
3 & 3 & $3{-}5$: 5 unvisited & tree edge, \texttt{par[5] = 3}; stack \texttt{1 3 5} \\
4 & 5 & $5{-}3$: parent & skip \\
5 & 5 & $5{-}1$: 1 visited, not the parent & \textbf{back edge}: cycle \texttt{1}, \texttt{5}, \texttt{par[5]=3}, \texttt{par[3]=1} \\
\bottomrule
\end{tabular}
\end{center}
Output \texttt{4} and \texttt{1 5 3 1}. Cities 2 and 4 are never visited: the search stops at the first cycle.
\end{tracebox}

\subsubsection*{6. Test suite}

\begin{testbox}
\centering\small
\begin{tabular}{@{}p{6.4cm}lp{5cm}@{}}
\toprule
\textbf{Input} & \textbf{Our output} & \textbf{Why} \\
\midrule
\texttt{2 1 / 1 2} & \texttt{IMPOSSIBLE} & one road: walking it back is not a round trip \\
\texttt{3 2 / 1 2 / 2 3} & \texttt{IMPOSSIBLE} & a path \\
\texttt{5 4 / 1 2 / 1 3 / 1 4 / 1 5} & \texttt{IMPOSSIBLE} & a star: many visited neighbours of 1, all children \\
\texttt{3 3 / 1 2 / 2 3 / 3 1} & \texttt{4 / 1 3 2 1} & smallest cycle \\
\texttt{6 5 / 1 2 / 2 3 / 4 5 / 5 6 / 6 4} & \texttt{4 / 4 6 5 4} & the cycle is not in vertex 1's component \\
\texttt{4 4 / 1 2 / 2 3 / 3 4 / 4 2} & \texttt{4 / 2 4 3 2} & the cycle does not contain the root \\
path $1..10^5$ + road $10^5{-}1$ & $k = 100001$ & deepest DFS, longest cycle \\
\bottomrule
\end{tabular}
\end{testbox}
All outputs come from running Approach~2. In addition, Approaches 1--3 were run on $3000$ random graphs with $n \le 9$; a checker verified every printed route (an edge between consecutive cities, no repeated city, at least 3 cities) and compared ``cycle or not'' with a DSU.

\subsubsection*{7. Common mistakes}

\begin{warningbox}[Mistake 1 --- no parent skip]
Without \texttt{if (v == par[u]) continue;}, the very first tree edge is reported: on the sample the program would print the ``round trip'' \texttt{1 3 1}, which visits only one other city.
\end{warningbox}

\begin{warningbox}[Mistake 2 --- searching only from city 1]
The fifth test above has its only cycle in the component $\{4,5,6\}$. A single DFS from 1 prints \texttt{IMPOSSIBLE}.
\end{warningbox}

\begin{warningbox}[Mistake 3 --- forgetting to close the route]
The judge wants $k$ cities \emph{including} the start twice: \texttt{4 / 1 5 3 1}, not \texttt{3 / 1 5 3}.
\end{warningbox}

\begin{warningbox}[Mistake 4 --- walking the parents past $v$]
The loop must stop at $v$ (\texttt{x != v}), not at the root. Walking to the root prints a ``cycle'' that leaves the cycle and comes back along a tree path, repeating cities.
\end{warningbox}

\subsubsection*{8. Takeaways}

\begin{itemize}[leftmargin=*]
    \item In an undirected DFS, a visited neighbour that is not the parent is an ancestor: the cycle is the tree path between them plus that edge.
    \item ``Two or more other cities'' is the statement's way of saying ``simple graph, ignore the parent edge''.
    \item Print by walking parents from $u$ up to $v$; close the loop for the output format.
    \item $n = 10^5$ means the DFS can be $10^5$ deep: iterative, or a large stack.
\end{itemize}

\begin{variantsbox}
\variant{One-way roads}{\varin{https://cses.fi/problemset/task/1678}}
  {every road is directed.}
  {the same search, but a visited vertex closes a cycle only if it is still on the stack.}
  {colours 0/1/2 and a gray target (Problem~\ref{ssec:ch2-round-trip-ii}).}
\variant{Parallel roads allowed}{\varnone}
  {two cities may be joined by several roads; a round trip may then use two of them.}
  {cycles of length 2 now exist.}
  {store edge ids and skip only the edge used to enter $u$ (template 1 of the pattern, which finds \texttt{6 7} on two parallel edges).}
\variant{Exactly one cycle, connected}{\varin{https://codeforces.com/problemset/problem/103/B}}
  {decide whether the graph is connected and has exactly one cycle.}
  {count instead of search: a connected graph has $n - 1$ edges plus one per independent cycle.}
  {connected and $m = n$ (Exercise E1.2).}
\variant{Components without a cycle}{\varin{https://codeforces.com/problemset/problem/659/E}}
  {count the components that are trees.}
  {per component, compare its edges and vertices.}
  {one DFS per component counting both; a tree has $e = v - 1$ (Exercise E1.4).}
\variant{A round trip through city $x$}{\varnone}
  {the trip must start and end in a given city $x$.}
  {$x$ needs two neighbours $a \ne b$ that stay connected when $x$ is removed.}
  {delete $x$, label the components of the rest, and look for two neighbours of $x$ with the same label; the cycle is $x, a \leadsto b, x$. \bigO{n+m}; checked against Approach~1 restricted to edges at $x$ on random graphs.}
\end{variantsbox}
